\section{Przegląd pisemnictwa}\label{Przegląd pisemnictwa}

W poniższym dziale znajduje się przegląd pisemnictwa z zakresu suplementacji jak i tematyk powiązanych.

%%%%%%%%%%%%%%%%%%%%%%%%%%%%%%%%%%%%%%%%%%%%%%%%%%
% I.1. Podstawy:
% - Definicje prawne suplementów diety i produktów leczniczych
% - Formy występowania suplementów oraz drogi dystrybucji
%%%%%%%%%%%%%%%%%%%%%%%%%%%%%%%%%%%%%%%%%%%%%%%%%%

\subsection{Suplementy diety}

Przed rozwinięciem tematyki pracy należy zwrócić uwagę na znaczące różnicę pomiędzy dwoma zbliżonymi do siebie kategoriami produktów występującymi na rynku które ze względu na liczne podobieństwa są często mylone przez konsumentów \cite{suplementy-a-leki}. Są nimi suplementy diety i produkty lecznicze. Oba mogą zawierać te same substancje, występować w tej samej formie i być sprzedawane w tych samych miejscach, jednak ich status prawny, przeznaczenie oraz sposób regulacji różnią się istotnie.

Na początek należy przywołać ich definicje prawne:
\begin{itemize}
    \item Suplement diety zgodnie z definicją zawartą w ustawie o bezpieczeństwie żywności i żywienia to "środek spożywczy, którego celem jest uzupełnienie normalnej diety, będący skoncentrowanym źródłem witamin lub składników mineralnych lub innych substancji wykazujących efekt odżywczy lub inny fizjologiczny, pojedynczych lub złożonych, wprowadzany do obrotu w formie umożliwiającej dawkowanie, w postaci: kapsułek, tabletek, drażetek i w innych podobnych postaciach, saszetek z proszkiem, ampułek z płynem, butelek z kroplomierzem i w innych podobnych postaciach płynów i proszków przeznaczonych do spożywania w małych, odmierzonych ilościach jednostkowych, z wyłączeniem produktów posiadających właściwości produktu leczniczego w rozumieniu przepisów prawa farmaceutycznego" \cite{definicja-suplementy-pl}.
    \item Product leczniczy (lek) według ustawy o prawie farmaceutycznym to "substancja lub mieszanina substancji, przedstawiana jako posiadająca właściwości zapobiegania lub leczenia chorób występujących u ludzi lub zwierząt lub podawana w celu postawienia diagnozy lub w celu przywrócenia, poprawienia lub modyfikacji fizjologicznych funkcji organizmu poprzez działanie farmakologiczne, immunologiczne lub metaboliczne" \cite{definicja-lek-pl}.
\end{itemize}

Do najważniejszych różnic pomiędzy oboma produktami należą:
\begin{itemize}
    \item Definicja prawna – lek podlega ustawie Prawo farmaceutyczne i posiada właściwości lecznicze, zapobiegawcze lub modyfikujące funkcje fizjologiczne; suplement diety jest definiowany w ustawie o bezpieczeństwie żywności i żywienia jako grupa żywności, czyli skoncentrowane źródło witamin, minerałów lub innych substancji o efekcie odżywczym.
    % W lekach dawka składników aktywnych wykazuje działanie farmakologiczne, terapeutyczne, natomiast w suplementach diety dawki składników mają charakter żywieniowy.
    \item Dawka i działanie – leki zazwyczaj zawierają wyższe dawki substancji czynnych i wykazują farmakologiczne, terapeutyczne działanie; suplementy diety podają dawki o charakterze żywieniowym i nie mogą sugerować leczenia ani zapobiegania chorobom.
    % Do leku natomiast dołączone jest oświadczenie medyczne, najczęściej w postaci ulotki w opakowaniu, które stwierdza, sugeruje lub daje do zrozumienia, że produkt lub jego składnik (składniki) mają własności leczenia lub zapobiegania chorobie (chorobom).
    \item Oświadczenia medyczne – leki dołączają ulotki, w których mogą stwierdzać właściwości lecznicze; suplementy diety nie mogą zawierać żadnych oświadczeń o leczeniu, zapobieganiu lub wyleczeniu chorób.
    % Znakowanie i sposób znakowania, prezentacja i reklama suplementów diety nie mogą przypisywać tym środkom właściwości zapobiegania, leczenia lub wyleczenia chorób człowieka, ani też sugerować takich właściwości
    \item Prezentacja i reklama – leki mogą używać języka medycznego i podkreślać wskazania terapeutyczne; suplementy diety muszą unikać stwierdzeń leczniczych i podkreślać jedynie swoje przeznaczenie jako uzupełnienie diety.
    \item Procedura rejestracji – jeżeli produkt spełnia kryteria leku, podlega wyłącznie przepisom dotyczącym produktów leczniczych i wymaga rejestracji w Urzędzie Rejestracji Leków; suplementy diety podlegają przepisom żywnościowym, wymagają zgłoszenia do Głównego Inspektora Sanitarnego, ale nie pełnej oceny klinicznej leku.
    % Składniki oraz ich poziomy zawarte w suplementach diety muszą zagwarantować, że normalne stosowanie tych produktów zgodnie ze wskazaniami producenta będzie bezpieczne dla konsumenta.
    \item Ryzyko i przeciwwskazania – leki podlegają ocenie ryzyka opartej na danych klinicznych, posiadają wyraźne przeciwwskazania i monitorowanie działań niepożądanych; suplementy diety mogą stwarzać ryzyko jedynie przy nadmiernym spożyciu (np. witamina A, żelazo) i ich przeciwwskazania oraz ostrzeżenia są mniej szczegółowo regulowane.
    \item Przeznaczenie – lek służy do leczenia, diagnozowania i modyfikacji funkcji organizmu; suplement diety ma na celu uzupełnienie diety, wsparcie odżywcze i ogólną poprawę stanu zdrowia.
\end{itemize}

%%%%%%%%%%%%%%%%%%%%%%%%%%%%%%%%%%%%%%%%%%%%%%%%%%
% I.2. Zasady suplementacji:
% - Racjonalna a nieracjonalna suplementacja
% - Trendy dietetyczne w społeczności (weganizm, wegetarianizm) a potrzeba suplementacji
%%%%%%%%%%%%%%%%%%%%%%%%%%%%%%%%%%%%%%%%%%%%%%%%%%

\subsection{Zasady suplementacji}

% Kiedy suplementacja jest uzasadniona?

Kluczowym pytaniem, które powinno poprzedzać każdą decyzję suplementacyjną, jest: czy mam niedobór lub czy jestem w grupie zwiększonego ryzyka niedoboru? Suplementacja bez wcześniej stwierdzonego deficytu lub wyraźnie zwiększonego zapotrzebowania jest w większości przypadków nieuzasadniona i może prowadzić do niepotrzebnego narażenia na ryzyko \cite{jarosz2020normy}.

W polskiej populacji kilka składników odżywczych ma rzeczywiste populacyjne wskazania do suplementacji. Witamina D jest tego najlepszym przykładem: ze względu na położenie geograficzne kraju, styl życia i dietę niedobór jest powszechny i uzasadnia suplementację od października do końca kwietnia, a u wielu osób przez cały rok \cite{pludowski2023guidelines}. Kwas foliowy jest rekomendowany kobietom planującym ciążę i w jej pierwszym trymestrze \cite{jarosz2020normy}. Witamina B12 powinna być suplementowana przez osoby na dietach wegańskich i wegetariańskich \cite{tardy2020vitamins}. Jod może wymagać suplementacji u osób unikających jodowanej soli i produktów mlecznych \cite{jarosz2020normy}. Żelazo jest wskazane wyłącznie przy potwierdzonym niedoborze lub niedokrwistości \cite{iolascon2024recommendations}.

Żywność powinna pozostawać priorytetowym źródłem składników odżywczych – dostarcza je w formach lepiej przyswajalnych i synergicznie działających z innymi składnikami produktu. Suplement może uzupełniać dietę, gdy ta jest rzeczywiście niedoborowa lub niemożliwa do zbilansowania z innych powodów. Nie powinien jej zastępować.

% Zasady bezpiecznego stosowania suplementów

Na podstawie dostępnych danych naukowych można sformułować kilka praktycznych zasad bezpiecznej suplementacji.

\begin{enumerate}
    \item Nie suplementuj bez wskazania. Jeśli nie masz potwierdzonego niedoboru ani zwiększonego zapotrzebowania, większość suplementów nie przyniesie ci korzyści, a może zaszkodzić \cite{jarosz2020normy,navarro2017liver}.
    \item Dobieraj formę i dawkę do rzeczywistych potrzeb. Magnez w postaci cytrynianu lub glicynianu wchłania się znacznie lepiej niż tlenek magnezu. Witamina D3 jest skuteczniejsza niż D2 \cite{pludowski2023guidelines,tardy2020vitamins}. Żelazo w formie bisglicynianu jest lepiej tolerowane niż siarczan żelazawy \cite{iolascon2024recommendations}.
    \item Informuj lekarza i farmaceutę o wszystkich stosowanych suplementach – szczególnie jeśli przyjmujesz leki na stałe.
    \item Nie przekraczaj zalecanych dawek, a najlepiej stosuj dawki zbliżone do RDA, nie do UL.
    \item Kupuj suplementy od sprawdzonych producentów, najlepiej w aptekach, nie przez niezweryfikowane sklepy internetowe – ryzyko fałszowania i zanieczyszczeń jest w tym kanale wyższe \cite{lin2023hidden}.
    \item Traktuj suplementację jako coś, co może być potrzebne przez określony czas, nie na zawsze – po wyrównaniu niedoboru warto wykonać kontrolne badania laboratoryjne i ocenić, czy dalsza suplementacja jest uzasadniona.
\end{enumerate}

% Rola dietetyka i innych specjalistów

Decyzja o suplementacji powinna w idealnej sytuacji poprzedzić konsultację ze specjalistą – lekarzem, farmaceutą lub dietetykiem. Dietetyk jest w tej trójce osobą najlepiej przygotowaną do oceny całościowego sposobu żywienia i identyfikacji rzeczywistych niedoborów na tle diety. Niestety, jak pokazują badania, konsultacje ze specjalistami są wśród studentów zdecydowanie rzadsze niż sięganie po informacje w Internecie \cite{vzevzelj2018prevalence,sicinska2022dietary}.

Farmaceuta, jako dostępny bez umówienia i bezpłatny doradca, pełni ważną rolę w pierwszej linii doradztwa suplementacyjnego. Jednak, jak pokazują dane z 2024 roku, ponad 60\% lekarzy pierwszego kontaktu i dietetyków przyznaje, że nie czuje się wystarczająco kompetentnych, żeby w pełni oceniać preparaty probiotyczne, a analogicznie może być z innymi suplementami \cite{hill2014expert}. Oznacza to, że kształcenie specjalistów w tym zakresie powinno być wzmacniane.

Na poziomie uczelni wyższych warto rozważyć włączenie zagadnień suplementacyjnych do programów kształcenia – nie tylko na kierunkach medycznych, ale też ogólnodostępnych kursach zdrowia publicznego. Doświadczenia z Michigan State University (kurs HNF 102 „Suplementy diety: dowody vs. hype") wskazują, że ustrukturyzowana edukacja w tym zakresie istotnie poprawia umiejętność krytycznej oceny twierdzeń producenckich i weryfikacji źródeł informacji \cite{becker2024development}.

%%%%%%%%%%%%%%%%%%%%%%%%%%%%%%%%%%%%%%%%%%%%%%%%%%
% I.3. Rodzaje suplementów
% - Omówienie korzyści i zagrożeń najczęściej stosowanych suplementów m.in. witamina D3, witaminy grupy B, witamina C, żelazo, magnez, cynk, kwasy omega 3, kreatyna, kolagen, probiotyki
%%%%%%%%%%%%%%%%%%%%%%%%%%%%%%%%%%%%%%%%%%%%%%%%%%

\subsection{Rodzaje suplementów}

%%%%%%%%%%%%%%%%%%%%%%%%%%%%%%%%%%%%%%%%%%%%%%%%%%
% I.4. Zagrożenia:
% - Ryzyko kumulacji dawek przy braniu różnych suplementów o nakładających się składach w szczególności substancje o wąskim zakresie terapeutycznym (witamina A, B6 oraz żelazo)
% - Interakcje suplementów z lekami
%%%%%%%%%%%%%%%%%%%%%%%%%%%%%%%%%%%%%%%%%%%%%%%%%%

\subsection{Zagrożenia}

%%%%%%%%%%%%%%%%%%%%%%%%%%%%%%%%%%%%%%%%%%%%%%%%%%
% I.5. Suplementacja u mężczyzn i kobiet (zainteresowanie zdrowiem i profilaktyką, marketing etc.)
%%%%%%%%%%%%%%%%%%%%%%%%%%%%%%%%%%%%%%%%%%%%%%%%%%

\subsection{Suplementacja u mężczyzn i kobiet}

%%%%%%%%%%%%%%%%%%%%%%%%%%%%%%%%%%%%%%%%%%%%%%%%%%
% I.6. Studenci a suplementacja (tendencje, najczęstsze źródła wiedzy, dezinformacja)
% - Częstość stosowania suplementów diety
% - Najczęściej stosowane suplementy diety
% - Powody stosowania suplementów diety
% - Źródła informacji o suplementach diety
% - Poziom wiedzy studentów na temat suplementów diety
%%%%%%%%%%%%%%%%%%%%%%%%%%%%%%%%%%%%%%%%%%%%%%%%%%

\subsection{Studenci a suplementacja}

\subsubsection{Częstość stosowania suplementów diety}

Dane z różnych krajów wskazują, że suplementacja jest wśród studentów zjawiskiem powszechnym, choć jej skala zależy od kontekstu kulturowego i środowiska akademickiego. W Chorwacji suplementy stosowało 30,5\% studentów, najczęściej preparaty witaminowo-mineralne. Co istotne, głównym źródłem informacji był Internet, natomiast konsultacje ze specjalistami należały do rzadkości \cite{croatia-knowledge-and-use}.
W Jordanii odsetek studentów stosujących suplementy był znacznie wyższy i wynosił 60,9\%. Dominowały preparaty jednoskładnikowe, a poziom wiedzy dotyczącej zasad bezpiecznej suplementacji oceniono jako niski \cite{knowledge-and-use}. Wyniki te sugerują, że popularność suplementów nie zawsze idzie w parze z odpowiednią świadomością żywieniową, co może sprzyjać nieprawidłowym praktykom.
Podobne tendencje zaobserwowano w Serbii, gdzie 55,7\% studentów deklarowało stosowanie suplementów, głównie witamin i minerałów. Co ciekawe, wyższy odsetek użytkowników odnotowano wśród studentów kierunków medycznych, co może wskazywać na większą dostępność informacji, ale niekoniecznie na bardziej krytyczne podejście \cite{serbian-undergraduate}.
W Polsce suplementy stosowało 41\% studentów Uniwersytetu Warszawskiego, a ich użycie było związane m.in. z płcią, aktywnością fizyczną i stylem życia \cite{diet-lifestyle}. Wyniki te wpisują się w ogólny trend rosnącej popularności suplementacji wśród młodych dorosłych.

\subsubsection{Najczęściej stosowane suplementy diety}

Najczęściej wybierane są preparaty witaminowe i mineralne — zarówno jedno-, jak i wieloskładnikowe. Dotyczy to również studentów kierunków medycznych i farmaceutycznych, którzy, mimo większego dostępu do wiedzy, nie różnią się znacząco od pozostałych studentów pod względem częstości suplementacji \cite{serbian-undergraduate}.
W wielu badaniach podkreśla się, że decyzje o suplementacji wynikają częściej z przekonań o ogólnym „wzmacniającym” działaniu suplementów niż z rzeczywistych potrzeb żywieniowych. Może to wskazywać na pewną lukę między wiedzą teoretyczną a praktyką \cite{knowledge-and-use,serbian-undergraduate}.

\subsubsection{Powody stosowania suplementów diety}

Najczęściej deklarowane powody suplementacji obejmują poprawę zdrowia, zapobieganie niedoborom, wzmocnienie odporności, zwiększenie energii oraz poprawę wyglądu skóry, włosów i paznokci \cite{knowledge-and-use,pharmacist-knowledge-and-use}. W badaniu przeprowadzonym w Bejrucie aż 68\% studentów stosowało suplementy, kierując się głównie poprawą samopoczucia i odporności \cite{pharmacist-knowledge-and-use}. Motywacje te mają więc w dużej mierze charakter subiektywny i nie zawsze wynikają z obiektywnej oceny stanu zdrowia.

\subsubsection{Źródła informacji o suplementach diety}

Najczęściej wskazywanym źródłem informacji pozostaje Internet — w tym media społecznościowe, blogi i reklamy. Konsultacje z lekarzami, dietetykami czy farmaceutami są zdecydowanie rzadsze \cite{croatia-knowledge-and-use,diet-lifestyle}. Może to sprzyjać utrwalaniu błędnych przekonań dotyczących suplementacji, zwłaszcza wśród osób o ograniczonej wiedzy żywieniowej.

\subsubsection{Poziom wiedzy na temat suplementów diety}

Mimo powszechnego stosowania suplementów poziom wiedzy studentów często okazuje się niewystarczający. Respondenci miewają trudności z odróżnieniem suplementów od leków, oceną bezpieczeństwa stosowania czy potencjalnych interakcji \cite{knowledge-and-use,serbian-undergraduate}.
Wśród studentów farmacji ponad połowa stosowała suplementy, jednak ich wiedza dotycząca dawkowania i działań niepożądanych była niepełna \cite{pharmacist-knowledge-and-use}. Podobne wnioski uzyskano w badaniu studentów Uniwersytetu Medycznego we Wrocławiu, którzy wskazywali na potrzebę dalszej edukacji w tym zakresie \cite{diet-lifestyle}. Wynika z tego, że nawet osoby kształcące się w obszarze zdrowia nie zawsze posiadają wystarczające kompetencje, by ocenić zasadność suplementacji.

\subsubsection{Uzasadnienie podjęcia tematu badawczego}

Przegląd literatury pokazuje, że suplementy diety są powszechnie stosowane przez studentów, jednak poziom ich wiedzy jest często niewystarczający, a głównym źródłem informacji pozostaje Internet. Taka sytuacja może sprzyjać nieprawidłowym praktykom suplementacyjnym, co uzasadnia potrzebę dalszych badań nad zależnościami między wiedzą, źródłami informacji a rzeczywistymi zachowaniami studentów.

